\documentclass[11pt,a4paper]{article}
\usepackage[utf8]{inputenc}
\usepackage[T1]{fontenc}
\usepackage{amsmath,amssymb,amsfonts,amsthm}
\usepackage{booktabs}
\usepackage{geometry}
\usepackage[colorlinks=true,linkcolor=blue,citecolor=blue,urlcolor=blue]{hyperref}
\geometry{margin=2.4cm}

\newtheorem{theorem}{Theorem}
\newtheorem{proposition}{Proposition}
\newtheorem{lemma}{Lemma}
\newtheorem{corollary}{Corollary}
\theoremstyle{definition}
\newtheorem{assumption}{Assumption}
\newtheorem{remark}{Remark}

\newcommand{\GeV}{\,\mathrm{GeV}}
\newcommand{\MeVtext}{\,\mathrm{MeV}}
\newcommand{\swsq}{\sin^2\theta_W}
\newcommand{\Minf}{M_\infty}
\newcommand{\Jvec}{\vec{J}}

\title{A Casimir Ladder for the Electroweak Bosons:\\
Exact Algebra of the de Vries--Rivero Ratio, a Boson Tower below $107\GeV$,\\
and Falsifiable Consequences}

\author{Claude Fable 5 (Anthropic)\thanks{This manuscript was produced by a
language model as a demonstration of autonomous theoretical synthesis, under
an explicit constraint of zero code execution: every closed form was derived
symbolically and every decimal obtained by hand arithmetic from those closed
forms. The source material is a corpus of research repositories maintained by
A.~Rivero; Section~\ref{sec:provenance} states which results are inherited,
which are reproved, and which are new. The bibliography was verified by a
second model (Claude Sonnet) against arXiv and the published record; an
adversarial referee report was solicited from an independent model (GPT-5.5,
via the Codex CLI) on a first draft, and this version incorporates the
resulting corrections --- including one genuine algebra error it caught. Both
reports are archived alongside the source.}}
\date{July 2, 2026}

\begin{document}
\maketitle

\begin{abstract}
We study the two-parameter algebraic structure underlying a family of
empirical electroweak mass relations first observed by H.~de~Vries and
A.~Rivero in 2004--2006 and since investigated, by several independent
routes, in a corpus of research programs. The structure is a one-line
kernel: to an internal label $j$ with Casimir $J=j(j+1)$ associate the
positive root $\lambda_+(J)$ of $\lambda^2+J\lambda-J=0$, and set
$m^2(j)=\Minf^2\,\lambda_+(J)$ for a single scale $\Minf$. With the channel
assignment $(\gamma,W,Z)=(0,\tfrac12,1)$ this gives
$m_W/m_Z=\sqrt{\lambda_+(3/4)/\lambda_+(2)}=0.881419\ldots$, within
$0.4\sigma$ of the 2024 on-shell world average. We prove three structural
results. (i)~The kernel is the unique attractor of a nested Feshbach
decimation with uniform rung data $(\sqrt J,\,J)$; we also record an
obstruction the source literature passes over --- a \emph{flat}
semi-infinite chain does not work, since the candidate boundary state is an
embedded resonance for both physical Casimirs. (ii)~The mixing parameter
implied by the kernel, $\swsq=1-\lambda_+(3/4)/\lambda_+(2)=
\bigl(19-3\sqrt{19}+3\sqrt{3}-\sqrt{57}\bigr)/16=0.2231013\ldots$, has
minimal polynomial $64x^4-304x^3+414x^2-190x+25$ over $\mathbb{Q}$: this is
precisely the quartic postulated independently in a second research line,
so two formulations believed merely numerically close are algebraically
identical; moreover the quartic's four roots are exactly the four branch
combinations $1-\lambda_a(J_W)/\lambda_b(J_Z)$, and the polynomial embeds
in an explicit two-parameter family whose coefficients are polynomial in
the channel Casimir and the ratio $\kappa=J_W/J_Z=3/8$ isolated
independently by a custodial-EFT analysis. (iii)~The negative branch obeys
two exact sum rules, and the
popular identifications of its levels with $v/2$ and $v/\sqrt2$ are
mutually inconsistent at the $1.8\%$ level, so at most one can be exact.
The positive branch predicts a tower of states at $96.5$, $99.6$, $101.5$,
$102.7\GeV$ accumulating at $\Minf=m_Z\sqrt{(1+\sqrt3)/2}=106.58\GeV$, with
a label-conservation pattern that suppresses single production of the
$j=3/2$ level --- compatible with the weak diphoton hints near
$95$--$96\GeV$ and decisively testable at FCC-ee, with LEP~2 recasts above
the $193\GeV$ pair threshold available already. We state which assumptions
remain underived, give the failure conditions, and enumerate the trials
factors any significance claim must price.
\end{abstract}

\section{Introduction}
\label{sec:intro}

The ratio of the $W$ and $Z$ boson masses is, in the Standard Model (SM), a
derived quantity: $m_W/m_Z=\cos\theta_W$ at tree level, with $\theta_W$
fixed by the ratio of two independent gauge couplings. The SM assigns it no
preferred value. It is therefore legitimate --- in the same exploratory
spirit in which Koide's charged-lepton relation has been studied for four
decades \cite{koide1982,koide1983prd,koide1990,koide2005,koide2018,foot1994}
--- to ask whether the measured value singles out a simple algebraic
structure, and, if one is found, to work out its full set of consequences so
that data can kill it or promote it.

An observation of this kind has an arXiv trail going back two decades. In
November 2004 H.~de~Vries noted \cite{devries2004} that the on-shell weak
mixing parameter agrees closely with a specific algebraic number built from
the two smallest nontrivial $SU(2)$ Casimirs; the observation was developed
in a joint preprint with A.~Rivero \cite{devriesrivero2005} and stated in
kernel form --- the eigenvalue problem used in this paper, with the value
$s^2_{dV}=0.22310132\ldots$ appearing verbatim --- in \cite{rivero2006}. A
corpus of research repositories has since attacked the structure from
several directions: as the secular equation of a two-channel kernel, as a
root of an explicit quartic polynomial, as an eigenvalue problem for a
Casimir-locked texture, and as a coefficient in a custodially symmetric
effective field theory. The connections between these formulations had been
checked numerically, never algebraically.

This paper does three things.

\emph{First}, it states the structure once, minimally, and derives its exact
spectrum in closed form (Sections~\ref{sec:kernel},~\ref{sec:spectrum}).
The kernel is a one-parameter family of $2\times2$ matrices; we prove it is
equivalently the unique attractor of a nested decimation recursion with
uniform rung data (Proposition~\ref{prop:ladder}), and we record an
obstruction that the source literature passes over: a flat semi-infinite
chain has the candidate state as an embedded resonance rather than an
eigenvalue (Remark~\ref{rem:notchain}), which constrains what a microscopic
origin can look like.

\emph{Second}, it proves the algebraic equivalence of the formulations
(Section~\ref{sec:theorem}). Theorem~\ref{thm:quartic} computes, via Galois
conjugation in the biquadratic field $\mathbb{Q}(\sqrt3,\sqrt{19})$, the
minimal polynomial of the mixing parameter implied by the kernel, and finds
exactly the quartic $64x^4-304x^3+414x^2-190x+25$ that a separate research
line had postulated as its fundamental object. What was believed to be two
compatible proposals is one proposal.

\emph{Third}, it extracts the falsifiable consequences derivable without new
dynamical input (Sections~\ref{sec:data}--\ref{sec:negative}): a
$2\times10^{-5}$-precision prediction for $m_W$ from $m_Z$ (at tree level,
on-shell --- the scheme caveat is discussed at length); a tower of states
between $96.5$ and $106.6\GeV$ with computable asymptotic spacing; a
label-conservation pattern suppressing their single production; two exact
sum rules for the negative branch; and a no-go
(Lemma~\ref{lem:incompatible}) showing that the two popular vacuum-scale
identifications of the negative branch cannot both be exact.

Section~\ref{sec:casimir} records a further coincidence: the coefficient
$\kappa=3/8=C_F/C_A$ of $SU(2)$, isolated independently by a custodial-EFT
analysis of the same data, is the ratio $J_W/J_Z$ of the very Casimirs that
enter the kernel. Section~\ref{sec:geometry} discusses which geometric
mechanisms could produce the kernel
\cite{manton1979,forgacsmanton1980,witten1981,salamstrathdee1982,duffnilssonpope1986,hosotani1983,rdss1983}
and states, as numbered assumptions, exactly what a derivation would have to
produce. Section~\ref{sec:falsification} collects the failure conditions,
and Section~\ref{sec:provenance} the provenance, method, and the trials
accounting that any significance claim must confront.

Throughout, ``prediction'' means: output of the closed formulas given here,
with stated inputs, at tree level, compared in the on-shell scheme. The
scheme question is the single largest interpretive uncertainty and is
treated in Section~\ref{sec:data}.

\section{The kernel and the decimation recursion}
\label{sec:kernel}

\subsection{Definition}

Fix a mass scale $\Minf$. To an internal label
$j\in\{0,\tfrac12,1,\tfrac32,\dots\}$ with Casimir $J=j(j+1)$ associate the
matrix
\begin{equation}
A(J)\;=\;\begin{pmatrix} 0 & \sqrt{J}\\[2pt] \sqrt{J} & -J\end{pmatrix},
\label{eq:kernel}
\end{equation}
acting on a two-dimensional space spanned by a ``visible'' level (first
slot) and a ``partner'' level (second slot), in units of $\Minf^2$. Its
characteristic equation is
\begin{equation}
\lambda^2 + J\lambda - J \;=\; 0,
\qquad
\lambda_\pm(J)\;=\;\frac{-J\pm\sqrt{J^2+4J}}{2},
\label{eq:secular}
\end{equation}
with $\lambda_+(J)\in(0,1)$ and $\lambda_-(J)<0$ for all $J>0$. The
physical mass-squared assignments studied here are
\begin{equation}
m^2(j) \;=\; \Minf^2\,\lambda_+\!\bigl(j(j+1)\bigr)
\qquad\text{(positive branch)},
\label{eq:masses}
\end{equation}
with the channel assignment
\begin{equation}
\gamma \leftrightarrow j=0,\qquad
W \leftrightarrow j=\tfrac12,\qquad
Z \leftrightarrow j=1,
\label{eq:channels}
\end{equation}
i.e.\ $J_\gamma=0$ (massless), $J_W=\tfrac34$, $J_Z=2$. The label $j$ is
distinct from the gauged weak isospin of the SM (under which $W^\pm,Z$ form
a triplet); it is a bookkeeping label of the construction, tentatively
interpreted as a representation label of an internal (custodial-like)
$SU(2)_i$, and assignment~\eqref{eq:channels} is an \emph{input}, inherited
from \cite{rivero2006}, whose derivation is an open problem. We return to
its status in Sections~\ref{sec:geometry} and~\ref{sec:provenance}.

Two elementary properties of \eqref{eq:secular} are used repeatedly:
\begin{equation}
\lambda_+ + \lambda_- = -J, \qquad \lambda_+\lambda_- = -J
\qquad\Longrightarrow\qquad
|\lambda_-| - \lambda_+ = J,\quad \lambda_+|\lambda_-| = J,
\label{eq:vieta}
\end{equation}
together with the rewriting $\lambda_+^2=J(1-\lambda_+)$ of
\eqref{eq:secular} itself.

\subsection{The resummation identity, and what it is not}

The matrix \eqref{eq:kernel} looks like a fine-tuned two-state system. The
following identity shows the same spectrum arises as the fixed point of a
self-similar decimation, which widens the space of candidate microscopic
origins.

\begin{proposition}[Resummation identity]
\label{prop:ladder}
Let $f(X)=J/(J+X)$ on $[0,\infty)$, for fixed $J>0$. Then:
\begin{enumerate}
\item[(i)] the fixed points of $f$ are exactly the roots
$\lambda_\pm(J)$ of \eqref{eq:secular};
\item[(ii)] for every $X_0\ge0$ the iterates $X_{k+1}=f(X_k)$ converge to
$\lambda_+(J)$;
\item[(iii)] $f$ is the static (zero-frequency) Schur reduction of the
two-level block
$\bigl(\begin{smallmatrix}0&\sqrt J\\ \sqrt J&-(J+X)\end{smallmatrix}\bigr)$
onto its visible slot: eliminating a partner that carries, besides the
Casimir mass $-J$, the effective mass $-X$ handed up from the stage below,
returns the effective visible mass $f(X)$. Hence any nested decimation
scheme with uniform rung data $(\sqrt J,\,J)$ has $\lambda_+(J)$ as its
unique attractor.
\end{enumerate}
\end{proposition}

\begin{proof}
(i) $f(X)=X$ is $X(J+X)=J$, i.e.\ \eqref{eq:secular}.
(ii) $f$ is a M\"obius transformation mapping $[0,\infty)$ into $(0,1]$,
strictly decreasing, with $f([0,1])=[J/(J{+}1),1]$. Hence $f\circ f$ is
increasing, so the even and odd subsequences of $(X_k)_{k\ge1}$ are
monotone and confined to $[J/(J{+}1),1]$; each converges to a fixed point
of $f\circ f$. Now $f\circ f$ is itself a M\"obius transformation and is
distinct from the identity (e.g.\ $f(f(0))=J/(J{+}1)\ne0$), so it has at
most two fixed points; the two fixed points $\lambda_\pm$ of $f$ already
account for them, leaving no room for genuine $2$-cycles. Both subsequence
limits lie in $(0,1]$ while $\lambda_-<0$, so both equal $\lambda_+$.
(For the record, the multiplier is
$|f'(\lambda_+)|=\lambda_+^2/J=1-\lambda_+<1$, using
$\lambda_+^2=J(1-\lambda_+)$.)
(iii) The Schur complement of the stated block onto the first slot at zero
frequency is $0-\sqrt J\,\bigl(-(J+X)\bigr)^{-1}\sqrt J=J/(J+X)=f(X)$.
\end{proof}

\begin{remark}[What the ladder is not]
\label{rem:notchain}
It is tempting to read Proposition~\ref{prop:ladder} as the statement that
a \emph{flat} semi-infinite chain --- the tridiagonal operator $T$ with
diagonal $(0,-J,-J,\dots)$ and off-diagonals $\sqrt J$ --- possesses
$\lambda_+$ as a boundary-localized eigenvalue. It does not. The tail of
$T$ has essential spectrum $[-J-2\sqrt J,\,-J+2\sqrt J\,]$, and
$\lambda_+$ lies inside this band precisely when
$(\lambda_++J)^2<4J$, which by $\lambda_++J=J/\lambda_+$ and
$\lambda_+^2=J(1-\lambda_+)$ reduces to $\lambda_+<\tfrac34$, i.e.\
$J<\tfrac94$. Both physical values $J_W=\tfrac34$
($\lambda_+=0.5687$) and $J_Z=2$ ($\lambda_+=0.7321$) satisfy this, so on
the flat chain the candidate boundary state is an embedded resonance
rather than an isolated eigenvalue. The correct reading of
Proposition~\ref{prop:ladder} is therefore recursive-hierarchical, as in
item (iii): each partner carries the effective mass of the stage below.
Microscopic proposals must implement a nested decimation --- successive
scales of a Kaluza--Klein-type recursion are the natural candidate ---
rather than a flat chain. We state this obstruction explicitly because the
source corpus passes over it.
\end{remark}

\section{Exact spectrum and the number field}
\label{sec:spectrum}

Evaluating \eqref{eq:secular} at the two physical Casimirs:
\begin{align}
\lambda_+\!\left(\tfrac34\right) &= \frac{\sqrt{57}-3}{8} = 0.568729\ldots,
&
\lambda_-\!\left(\tfrac34\right) &= -\frac{3+\sqrt{57}}{8} = -1.318729\ldots,
\label{eq:lam34}\\[2pt]
\lambda_+(2) &= \sqrt{3}-1 = 0.732051\ldots,
&
\lambda_-(2) &= -(1+\sqrt{3}) = -2.732051\ldots
\label{eq:lam2}
\end{align}
The mass-squared ratio on the positive branch is
\begin{equation}
\rho \;\equiv\; \frac{m_W^2}{m_Z^2}
\;=\;\frac{\lambda_+(3/4)}{\lambda_+(2)}
\;=\;\frac{\sqrt{57}-3}{8(\sqrt3-1)}
\;=\;\frac{3\sqrt{19}+\sqrt{57}-3\sqrt3-3}{16}
\;=\;0.776899\ldots,
\label{eq:rho}
\end{equation}
where we rationalized with $(\sqrt3-1)^{-1}=(\sqrt3+1)/2$ and used
$\sqrt{57}=\sqrt3\sqrt{19}$. Hence
\begin{equation}
\frac{m_W}{m_Z} \;=\;\sqrt{\rho}\;=\;0.8814186\ldots,
\qquad
\swsq \;\equiv\; 1-\rho
\;=\;\frac{19-3\sqrt{19}+3\sqrt3-\sqrt{57}}{16}
\;=\;0.2231013\ldots
\label{eq:sw}
\end{equation}
All these numbers live in the real biquadratic field
$K=\mathbb{Q}(\sqrt3,\sqrt{19})$, of degree $4$ over $\mathbb{Q}$, with
Galois group $\mathrm{Gal}(K/\mathbb{Q})\cong
\mathbb{Z}/2\times\mathbb{Z}/2$ generated by
\[
\sigma:\ \sqrt3\mapsto-\sqrt3,\qquad
\tau:\ \sqrt{19}\mapsto-\sqrt{19},
\]
each of which flips $\sqrt{57}$. A curiosity worth recording: since $K$ is
biquadratic, every quantity in this construction is expressible by nested
square roots alone --- in this model $\swsq$ is a ruler-and-compass
constructible number.

\section{The equivalence theorem}
\label{sec:theorem}

A second research line in the corpus posits, as its fundamental object, the
quartic polynomial
\begin{equation}
P(x) \;=\; 64x^4-304x^3+414x^2-190x+25,
\label{eq:quartic}
\end{equation}
one of whose roots is matched against electroweak data. The two lines were
believed compatible because their decimal expansions agree. We now prove
they are the same object.

\begin{theorem}[The quartic is the minimal polynomial of the mixing
parameter]
\label{thm:quartic}
Let $s=\swsq=\bigl(19-3\sqrt{19}+3\sqrt3-\sqrt{57}\bigr)/16$ as in
\eqref{eq:sw}. Then the minimal polynomial of $s$ over $\mathbb{Q}$ is
exactly $P(x)$ of \eqref{eq:quartic}, normalized to leading coefficient
$64$. Consequently the kernel formulation \eqref{eq:kernel} with channels
\eqref{eq:channels} and the quartic formulation \eqref{eq:quartic} define
the same algebraic number, and the two research lines study one structure.
\end{theorem}

\begin{proof}
The four Galois images of $s$ under
$\mathrm{Gal}(K/\mathbb{Q})=\{1,\sigma,\tau,\sigma\tau\}$ are obtained by
the sign choices $\varepsilon_3,\varepsilon_{19}\in\{\pm1\}$ on $\sqrt3$
and $\sqrt{19}$ (the sign of $\sqrt{57}$ being
$\varepsilon_3\varepsilon_{19}$):
\[
s(\varepsilon_3,\varepsilon_{19})
=\frac{19-3\varepsilon_{19}\sqrt{19}+3\varepsilon_3\sqrt3
-\varepsilon_3\varepsilon_{19}\sqrt{57}}{16}.
\]
Write $s_1=s(+,+)=s$, $s_2=s(-,+)$, $s_3=s(+,-)$, $s_4=s(-,-)$.

\emph{The orbit has size four.} Since $\{1,\sqrt3,\sqrt{19},\sqrt{57}\}$ is
a $\mathbb{Q}$-basis of $K$, two images $s(\varepsilon)$,
$s(\varepsilon')$ coincide only if their coefficient vectors coincide,
i.e.\ only if $\varepsilon=\varepsilon'$; the four images are pairwise
distinct. (Numerically: $s_1=0.223101$, $s_2=0.517311$, $s_4=1.208169$,
$s_3=2.801418$.) Hence the stabilizer of $s$ in the Galois group is
trivial, $[\mathbb{Q}(s):\mathbb{Q}]=4$, and the degree-$4$ monic
polynomial $\prod_i(x-s_i)$, which is Galois-invariant and therefore has
rational coefficients, is irreducible over $\mathbb{Q}$ and is the minimal
polynomial of $s$.

\emph{Computation of the coefficients.} We compute the elementary symmetric
functions by grouping conjugate pairs.

Sum: the $\sqrt3,\sqrt{19},\sqrt{57}$ terms cancel over the four sign
patterns, so $e_1=4\cdot 19/16=19/4$.

Pair data at fixed $\varepsilon_{19}$: with
$A=19-3\varepsilon_{19}\sqrt{19}$ and
$B=3\sqrt3-\varepsilon_{19}\sqrt{57}$, the two values of $\varepsilon_3$
give $(A\pm B)/16$, so
\[
s(+,\varepsilon_{19})\,s(-,\varepsilon_{19})=\frac{A^2-B^2}{256},
\qquad
A^2=532-114\,\varepsilon_{19}\sqrt{19},
\qquad
B^2=84-18\,\varepsilon_{19}\sqrt{19},
\]
using $(3\sqrt3)^2=27$, $(\sqrt{57})^2=57$ and
$2\cdot3\sqrt3\cdot\sqrt{57}=6\sqrt{171}=18\sqrt{19}$. Thus
\[
s_1s_2=\frac{28-6\sqrt{19}}{16},\qquad
s_3s_4=\frac{28+6\sqrt{19}}{16},\qquad
s_1+s_2=\frac{19-3\sqrt{19}}{8},\qquad
s_3+s_4=\frac{19+3\sqrt{19}}{8}.
\]

Constant term:
$e_4=s_1s_2\,s_3s_4=\dfrac{28^2-36\cdot19}{256}
=\dfrac{784-684}{256}=\dfrac{25}{64}$.

Second symmetric function:
\[
e_2=s_1s_2+s_3s_4+(s_1+s_2)(s_3+s_4)
=\frac{56}{16}+\frac{361-171}{64}
=\frac{7}{2}+\frac{190}{64}=\frac{207}{32}.
\]

Third symmetric function:
\[
e_3=s_1s_2(s_3+s_4)+s_3s_4(s_1+s_2)
=\frac{(28-6\sqrt{19})(19+3\sqrt{19})+(28+6\sqrt{19})(19-3\sqrt{19})}{128}
=\frac{(190-30\sqrt{19})+(190+30\sqrt{19})}{128}=\frac{95}{32}.
\]

Therefore the minimal polynomial is
\[
x^4-\frac{19}{4}x^3+\frac{207}{32}x^2-\frac{95}{32}x+\frac{25}{64}
\;=\;\frac{1}{64}\,P(x). \qedhere
\]
\end{proof}

\begin{corollary}
\label{cor:rho}
(a) The mass-squared ratio $\rho=1-s$ has minimal polynomial
$64y^4+48y^3-114y^2+18y+9$, obtained by expanding $P(1-y)$.
(b) The mass ratio $z=m_W/m_Z=\sqrt{\rho}$ has degree $8$ over
$\mathbb{Q}$, with minimal polynomial $64z^8+48z^6-114z^4+18z^2+9$.
\end{corollary}

\begin{proof}
(a) Direct substitution; the coefficients follow from the binomial
expansions of $(1-y)^k$, $k\le4$.
(b) $z^2=\rho$ with $[\mathbb{Q}(\rho):\mathbb{Q}]=4$, so
$[\mathbb{Q}(z):\mathbb{Q}]\in\{4,8\}$, and it is $8$ unless $\rho$ is a
square in $K=\mathbb{Q}(\rho)$. Suppose $\rho=\beta^2$ with $\beta\in K$.
$K$ is a subfield of $\mathbb{R}$ and Galois over $\mathbb{Q}$, so every
Galois image of $\rho$ would be the square of a real number, hence
nonnegative. But the image of $\rho=1-s$ under $\tau$ is
$1-s_3=-1.801418\ldots<0$. Contradiction; so
$[\mathbb{Q}(z):\mathbb{Q}]=8$ and the stated monic-normalized degree-$8$
polynomial, of which $z$ is a root, is its minimal polynomial.
\end{proof}

\begin{remark}
Theorem~\ref{thm:quartic} is a statement about the internal consistency of
the theoretical corpus, and it \emph{removes} one apparent independent
confirmation: the quartic line and the kernel line are one coincidence, so
the naive over-determination count goes down even as the structural economy
goes up. The bookkeeping consequences are drawn in
Section~\ref{sec:provenance}.
\end{remark}

\subsection{The branch structure of the quartic: a two-parameter family}
\label{sec:branchquartic}

The proof of Theorem~\ref{thm:quartic} contains more information than the
statement. Tracking the Galois action on the closed forms
\eqref{eq:lam34}--\eqref{eq:lam2}: $\tau$ (which flips $\sqrt{57}$ but fixes
$\sqrt3$) swaps $\lambda_+(\tfrac34)\leftrightarrow\lambda_-(\tfrac34)$
while fixing $\lambda_\pm(2)$; $\sigma\tau$ swaps
$\lambda_+(2)\leftrightarrow\lambda_-(2)$ while fixing
$\lambda_\pm(\tfrac34)$; and $\sigma$ swaps both. Hence the four roots of
the corpus quartic are precisely the four \emph{branch combinations}
\[
s_{ab}\;=\;1-\frac{\lambda_a(J_W)}{\lambda_b(J_Z)},
\qquad a,b\in\{+,-\}:
\]
$s_{++}=\swsq=0.223101$, $\;s_{--}=0.517311$, $\;s_{+-}=1.208169$,
$\;s_{-+}=2.801418$, and the Galois group of Theorem~\ref{thm:quartic} acts
as the group of branch swaps. The ``spurious'' roots of the quartic are
the other branch choices --- the arithmetic of the polynomial is the
physics of the branch structure. This observation globalizes:

\begin{theorem}[Two-parameter branch quartic]
\label{thm:family}
Fix channel Casimirs $I,K>0$ and let $u_\pm=\lambda_\pm(I)$,
$w_\pm=\lambda_\pm(K)$ be the roots of $p_I(x)=x^2+Ix-I$ and
$p_K(x)=x^2+Kx-K$. Set $\kappa=I/K$ and $s_{ab}=1-u_a/w_b$. Then,
identically in $I,K$,
\[
\prod_{a,b\in\{\pm\}}\bigl(x-s_{ab}\bigr)
\;=\;
x^4-(4+I)\,x^3
+\Bigl[\,6+2I-\kappa\,(I+2)\Bigr]x^2
-(4+I)(1-\kappa)\,x
+(1-\kappa)^2 .
\]
At $(I,K)=(J_W,J_Z)=(\tfrac34,2)$, i.e.\ $\kappa=\tfrac38$, this is
$P(x)/64$ with $P$ the corpus quartic \eqref{eq:quartic}.
\end{theorem}

\begin{proof}
Vieta for the two channels gives $u_++u_-=-I$, $u_+u_-=-I$,
$w_++w_-=-K$, $w_+w_-=-K$; hence the two identities used throughout:
$\tfrac1{w_+}+\tfrac1{w_-}=\tfrac{w_++w_-}{w_+w_-}=1$ and
$\tfrac1{w_+^2}+\tfrac1{w_-^2}=1+\tfrac2K$, together with
$u_+^2+u_-^2=I^2+2I$. The key on-shell identity: if $w^2+Kw-K=0$ then
$w^2=K(1-w)$, so
\[
p_I(w)=w^2+Iw-I=K-Kw+Iw-I=(K-I)(1-w).
\]
\emph{Constant term.} Since
$\prod_a(1-u_a/w)=\bigl[w^2-(u_++u_-)w+u_+u_-\bigr]/w^2=p_I(w)/w^2$,
\[
e_4=\prod_b\frac{p_I(w_b)}{w_b^2}
=\frac{(K-I)^2\prod_b(1-w_b)}{(w_+w_-)^2}
=\frac{(K-I)^2\,p_K(1)}{K^2}
=\Bigl(\frac{K-I}{K}\Bigr)^{\!2}=(1-\kappa)^2,
\]
using $p_K(1)=1+K-K=1$.
\emph{Sum.} $e_1=\sum_{ab}s_{ab}
=4-\bigl(\sum_a u_a\bigr)\bigl(\sum_b 1/w_b\bigr)=4-(-I)(1)=4+I$.
\emph{Second symmetric function.} From
$\sum_{ab}s_{ab}^2
=4-2\bigl(\sum u\bigr)\bigl(\sum 1/w\bigr)
+\bigl(\sum u^2\bigr)\bigl(\sum 1/w^2\bigr)
=4+2I+(I^2+2I)\bigl(1+\tfrac2K\bigr)$
and $e_2=\tfrac12\bigl(e_1^2-\sum s^2\bigr)$ one finds
$e_2=6+2I-(I^2+2I)/K=6+2I-\kappa(I+2)$.
\emph{Third symmetric function.} Using $e_3=e_4\sum_{ab}1/s_{ab}$ and
$1/s_{ab}=w_b/(w_b-u_a)$: for fixed $w=w_b$,
\[
\sum_a\frac{w}{w-u_a}=\frac{w\,p_I'(w)}{p_I(w)}
=\frac{w(2w+I)}{(K-I)(1-w)}
=\frac{(2K-I)(1-w)+I}{(K-I)(1-w)},
\]
where the numerator was rewritten on-shell:
$w(2w+I)=2w^2+Iw=2K(1-w)+Iw=2K-(2K-I)w=(2K-I)(1-w)+I$. Summing over $b$
with
$\sum_b 1/(1-w_b)=p_K'(1)/p_K(1)=2+K$ gives
\[
\sum_{ab}\frac1{s_{ab}}
=\frac{2(2K-I)+I(2+K)}{K-I}
=\frac{K(4+I)}{K-I},
\qquad
e_3=(1-\kappa)^2\cdot\frac{K(4+I)}{K-I}=(4+I)(1-\kappa).
\]
The four coefficients assemble into the stated quartic. Specializing
$I=\tfrac34$, $\kappa=\tfrac38$: $e_1=\tfrac{19}4$,
$e_2=6+\tfrac32-\tfrac38\cdot\tfrac{11}4=\tfrac{207}{32}$,
$e_3=\tfrac{19}4\cdot\tfrac58=\tfrac{95}{32}$,
$e_4=\tfrac{25}{64}$, matching Theorem~\ref{thm:quartic}. (As an
independent check of the general identity, at the degenerate point $I=K=2$
the formula gives roots $\{0,0,3-\sqrt3,3+\sqrt3\}$ with
$e_1=6$, $e_2=6$, $e_3=e_4=0$, which is what the branch values
$1-\lambda_a/\lambda_b$ of a single channel indeed are.)
\end{proof}

\begin{remark}
\label{rem:family}
Theorem~\ref{thm:family} demystifies the corpus quartic twice over. First,
the integers $(64,304,414,190,25)$ carry no independent information: they
are the image of the two-parameter family at $(I,\kappa)=(\tfrac34,
\tfrac38)$, with $64$ the least common denominator. Second, the
coefficients depend on the pair $(I,K)$ only through $I$ and the ratio
$\kappa=I/K$ --- the same $\kappa=\tfrac38$ isolated by the custodial-EFT
line (Section~\ref{sec:casimir}), which therefore sits literally inside
the quartic's $x^1$ and $x^0$ coefficients, $e_3=e_1(1-\kappa)$ and
$e_4=(1-\kappa)^2$. For the statistics program the family is equally
useful: it \emph{is} the natural pre-registerable hypothesis class (scan
$(I,K)$ over Casimir pairs) inside which the look-elsewhere cost of the
choice $(\tfrac34,2)$ can be priced without ambiguity.
\end{remark}

\section{Confrontation with data}
\label{sec:data}

The construction has one dimensionful input and, at tree level, one output
ratio; the on-shell relation $\swsq=1-m_W^2/m_Z^2$ makes the $m_W$
comparison and the on-shell-angle comparison the \emph{same} test, and we
present it once, as a prediction for $m_W$ given $m_Z$:
\begin{equation}
\Minf=\frac{m_Z}{\sqrt{\lambda_+(2)}}=m_Z\sqrt{\frac{1+\sqrt3}{2}},
\qquad
m_W^{\rm pred}=m_Z\,\sqrt{\rho}.
\label{eq:anchor}
\end{equation}
With $m_Z=91.1876(21)\GeV$ \cite{pdg2024},
\begin{equation}
\Minf = 106.578\GeV,
\qquad
m_W^{\rm pred} = 80.3745(19)\GeV,
\label{eq:mwpred}
\end{equation}
where the quoted uncertainty is input propagation only
($\partial m_W/\partial m_Z=\sqrt\rho$). Against the measured values (pull
$\equiv(m_W^{\rm pred}-m_W^{\rm meas})/\sigma_{\rm comb}$, combining the
experimental error with the $\pm0.0019$ input error in quadrature):

\begin{center}
\begin{tabular}{lccc}
\toprule
measurement & $m_W$ [GeV] & source & pull \\
\midrule
PDG 2024 world average & $80.3692(133)$ & \cite{pdg2024} & $+0.4\sigma$\\
CMS 2024 & $80.3602(99)$ & \cite{cmsw2024} & $+1.4\sigma$\\
PDG-live 2026 snapshot & $80.3625(77)$ & [corpus pin, unverified] & $+1.5\sigma$\\
CDF II 2022 & $80.4335(94)$ & \cite{cdf2022} & $-6.1\sigma$\\
\bottomrule
\end{tabular}
\end{center}

Equivalently, the predicted on-shell angle is $\swsq=0.2231013$, to be
compared with $0.22305(23)$ \cite{pdg2024} ($+0.2\sigma$); the corpus also
pins an unverified 2026 snapshot $0.223339(153)$ ($-1.6\sigma$). Three
honesty notes are mandatory.

\emph{(a) Scheme.} The algebraic prediction is a pure number; the data are
scheme- and scale-dependent. The comparison above uses the on-shell
definition, for which $\swsq\equiv1-m_W^2/m_Z^2$ is an identity. The
effective leptonic angle $\sin^2\theta^{\ell}_{\rm eff}\simeq0.2315$
differs from the on-shell value by known radiative corrections of order
$4\%$; even the mass \emph{definition} matters at this precision, since the
conventional Breit--Wigner (running-width) masses and the complex-pole
masses differ by $\Gamma^2/2m$, roughly $34\MeVtext$ for the $Z$ and
$27\MeVtext$ for the $W$ --- an order of magnitude above the experimental
errors --- and the
construction has not declared which definition it predicts. A claim that
the kernel value targets one scheme and one mass definition rather than
another is dynamical input the construction does not yet possess. There is
also a sharper internal tension: the SM global electroweak fit returns an
indirect $m_W\simeq80.355(6)\GeV$ \cite{pdg2024} from $m_Z$, $m_t$, $m_H$
and the couplings via the radiative $\Delta r$ relation, and the algebraic
value $80.3745$ lies about $3\sigma$ \emph{above} it; if the on-shell
prediction were exact, consistency would require a compensating
new-physics contribution to $\Delta r$ --- an unpleasant but usefully
concrete corollary. The necessary next step is a one-loop matching
computation determining whether any consistent scheme-and-definition
assignment survives at current precision.

\emph{(b) Data motion.} The pull moved from $+0.4\sigma$ (2024 world
average) to about $+1.5\sigma$ (CMS 2024 and the 2026 snapshot) while CDF~II
sits $6\sigma$ on the other side; the experimental situation for $m_W$ is
unsettled. A retrodiction study across PDG editions --- checking whether
the agreement improved as errors shrank, the signature of a real constraint
rather than a fit to current centrals --- is defined in the corpus's
planning documents and should be executed before stronger claims are made.

\emph{(c) Error budget.} The theory error (scheme choice, radiative
corrections, and the unknown dynamics behind the kernel) dominates the
$\pm0.0019$ input error and is presently unquantified. Any use of
\eqref{eq:mwpred} as a precision claim without that budget would be
overreach.

\section{The positive-branch tower}
\label{sec:tower}

If assignment \eqref{eq:channels} extends to higher $j$, the positive
branch is a tower of levels below $\Minf$:
\begin{equation}
m(j)=\Minf\sqrt{\lambda_+\bigl(j(j+1)\bigr)}\,,
\qquad
\lambda_+(J) = 1-\frac1J+\mathcal{O}\!\left(\frac1{J^2}\right),
\qquad
m(j)\;\simeq\;\Minf\Bigl(1-\frac{1}{2J}\Bigr)\ \ (J\to\infty).
\label{eq:towerasym}
\end{equation}
Exact closed forms and hand-evaluated values:
\begin{center}
\begin{tabular}{ccll}
\toprule
$j$ & $J$ & $\lambda_+(J)$ & $m(j)$ \\
\midrule
$\tfrac12$ & $\tfrac34$ & $(\sqrt{57}-3)/8=0.568729$ & $80.374\GeV$ ($=m_W^{\rm pred}$)\\
$1$ & $2$ & $\sqrt3-1=0.732051$ & $91.188\GeV$ ($=m_Z$, input)\\
$\tfrac32$ & $\tfrac{15}4$ & $(\sqrt{465}-15)/8=0.820482$ & $96.54\GeV$\\
$2$ & $6$ & $(\sqrt{60}-6)/2=0.872983$ & $99.58\GeV$\\
$\tfrac52$ & $\tfrac{35}4$ & $(\sqrt{1785}-35)/8=0.906157$ & $101.45\GeV$\\
$3$ & $12$ & $(\sqrt{192}-12)/2=0.928203$ & $102.68\GeV$\\
$\vdots$ & & & accumulating at $\Minf=106.58\GeV$\\
\bottomrule
\end{tabular}
\end{center}

The accumulation point carries a consistency obligation the source corpus
does not discuss: taken literally, \eqref{eq:towerasym} places infinitely
many levels below $107\GeV$, with spacing $\sim\Minf/2J$. Compatibility
with LEP and electroweak precision data then \emph{requires} the couplings
of the levels to SM fields to fall rapidly with $j$ --- which is exactly
what the label-conservation pattern of the next subsection supplies,
since reaching a level of internal label $j$ from SM states ($j\le1$)
demands increasingly label-violating operators. Either the tower
terminates dynamically, or its visibility must decay with $j$; a
derivation must produce one of the two, and precision $Z$-pole and LEP~2
data already bound the alternatives.

\subsection{A label-conservation pattern and the $95$--$96\GeV$ hints}
\label{sec:selection}

The following discussion is conditional on Assumption~\ref{ass:su2i} below
(approximate conservation of the internal label by the production and decay
dynamics) and is qualitative: the 4D spin, CP, and width of the tower
states are unspecified by the kernel, and rate predictions require the
label-breaking dynamics, which the construction does not yet possess. What
the label bookkeeping alone gives is a \emph{pattern}.

With the SM bosons carrying $j_\gamma=0$, $j_W=\tfrac12$, $j_Z=1$,
two-boson channels carry
\[
\tfrac12\otimes\tfrac12=0\oplus1,\qquad
1\otimes1=0\oplus1\oplus2,\qquad
\tfrac12\otimes1=\tfrac12\oplus\tfrac32,\qquad
1\otimes0=1,\qquad \tfrac12\otimes 0=\tfrac12 .
\]
Consequences for the electrically neutral tower members:
\begin{itemize}
\item The $j=\tfrac32$ level at $96.54\GeV$ appears in no neutral two-boson
channel: $\tfrac32\not\subset0\oplus1$ ($WW$, $Z\gamma$, $\gamma\gamma$)
and $\tfrac32\not\subset0\oplus1\oplus2$ ($ZZ$). The only two-boson channel
containing $\tfrac32$ is $W^\pm Z$, which is electrically charged and
available only to its charged partners. Label-conserving single production
of the neutral $j=\tfrac32$ state in Drell--Yan or vector-boson fusion is
therefore forbidden; it proceeds only through label breaking or
higher-dimensional operators, and $\gamma\gamma$ couplings arise at loop
level. A narrow, weakly produced state near $96.5\GeV$ is \emph{compatible}
with --- though in no way established by --- the unresolved low-mass
diphoton hints near $95$--$96\GeV$ (CMS \cite{cms95}: local $2.9\sigma$ at
$95.4\GeV$; ATLAS \cite{atlas95}: local $1.7\sigma$ at the same mass). The
pattern statement is falsifiable in shape: within this construction such a
state should remain weak in single production, while pair production opens
at $\sqrt{s}>2m(3/2)=193\GeV$.
\item The $j=2$ level at $99.58\GeV$ does occur in $1\otimes1$: it is the
first tower state with a label-allowed two-$Z$ vertex, so it would appear
in $ZZ^{*}$ and near-threshold $ZZ$ observables rather than in
$\gamma\gamma$.
\end{itemize}
Two data-driven tests follow. LEP~2 accumulated data up to
$\sqrt s=209\GeV$, above the $j=3/2$ pair threshold of $193\GeV$: a
dedicated recast of LEP~2 cross sections for pair-produced narrow neutral
states in the $96$--$104\GeV$ range would already constrain the tower and
is, to our knowledge, the sharpest retrodictive test available. FCC-ee,
with runs at $240\GeV$ and precision scans, either finds the tower or
excludes it up to the accumulation point $106.58\GeV$.

\section{The negative branch: sum rules and a no-go}
\label{sec:negative}

The corpus matches the negative-branch levels to the electroweak vacuum
expectation value $v=246.220\GeV$:
\begin{equation}
m_-(j)\equiv\Minf\sqrt{|\lambda_-(j(j+1))|}:\qquad
m_-(\tfrac12) = 122.39\GeV \ \leftrightarrow\ \tfrac v2=123.11\GeV,\quad
m_-(1) = 176.16\GeV \ \leftrightarrow\ \tfrac v{\sqrt2}=174.10\GeV,
\label{eq:negmatch}
\end{equation}
with deviations $-0.58\%$ and $+1.18\%$ respectively. The $J=2$
identification has a strikingly compact exact form worth recording: since
$|\lambda_-(2)|/\lambda_+(2)=(1+\sqrt3)/(\sqrt3-1)=(1+\sqrt3)^2/2$, it is
equivalent to
\begin{equation}
v \;=\; \bigl(1+\sqrt3\bigr)\,m_Z \;=\; |\lambda_-(2)|\; m_Z
\;=\;249.13\GeV \quad(\text{vs.\ }246.22,\ +1.18\%),
\label{eq:vformula}
\end{equation}
the vacuum expectation value as the negative-branch \emph{eigenvalue}
itself in units of $m_Z$. We emphasize the epistemic status: pending a
mechanism that interprets the negative branch, the identifications
\eqref{eq:negmatch} are numerology; what this section contributes is the
exact algebra they must satisfy, and a no-go they currently violate.

From \eqref{eq:vieta}, with no further input:
\begin{proposition}[Branch sum rules]
\label{prop:sumrules}
For every $j$, with $J=j(j+1)$,
\[
m_-^2(j)-m_+^2(j) \;=\; J\,\Minf^2,
\qquad\qquad
m_+(j)\,m_-(j) \;=\; \sqrt{J}\,\Minf^2 .
\]
\end{proposition}
\noindent
Inserting the identifications \eqref{eq:negmatch} and the measured
$m_W,m_Z$, the first rule holds at $+2.1\%$ ($J=\tfrac34$) and $-3.2\%$
($J=2$); the second at $+0.6\%$ and $-1.2\%$ --- i.e.\ exactly at the level
of the mass matches themselves, as must happen algebraically.

More constraining is that the two identifications fight each other:

\begin{lemma}[The two vacuum identifications cannot both be exact]
\label{lem:incompatible}
If $m_-(\tfrac12)=v/2$ and $m_-(1)=v/\sqrt2$ held simultaneously and
exactly, then
\[
\frac{|\lambda_-(2)|}{|\lambda_-(3/4)|}
=\frac{(v/\sqrt2)^2}{(v/2)^2}=2 .
\]
But the exact value is
\[
\frac{|\lambda_-(2)|}{|\lambda_-(3/4)|}
=\frac{8(1+\sqrt3)}{3+\sqrt{57}}
=2.07173\ldots,
\]
which differs from $2$ by $3.6\%$ (i.e.\ $1.8\%$ in the mass ratio ---
realized in the data as the observed $-0.58\%/+1.18\%$ split). Hence at
most one of the two identifications can be exact; the pair is either
approximate (radiatively corrected) or partially accidental.
\end{lemma}

This no-go, which to our knowledge is stated here for the first time,
sharpens the negative-branch question considerably: any future derivation
must either pick one exact identification and explain the other's
percent-level offset, or predict the specific combination
$8(1+\sqrt3)/(3+\sqrt{57})$ in place of $2$ --- e.g.\ the conventional
ratio $(v/\sqrt2)/(v/2)=\sqrt2=1.41421$ would be replaced by
$\sqrt{2.07173}=1.43935$.

\section{The Casimir coincidence \texorpdfstring{$\kappa=3/8$}{k=3/8}}
\label{sec:casimir}

A custodial-EFT analysis in the corpus, working from the same electroweak
data with entirely different methods (Wilson-coefficient fits plus a
look-elsewhere analysis), isolated the coefficient
\[
\kappa=\frac38=\frac{C_F}{C_A}\Big|_{SU(2)}
\]
as the group-theoretic content of its best-fit relation. We record that
this is the same pair of Casimirs that powers the kernel:
\[
\frac{J_W}{J_Z}=\frac{C_2(\text{fund})}{C_2(\text{adj})}
=\frac{3/4}{2}=\frac38 .
\]
Theorem~\ref{thm:family} makes the cross-link exact rather than
suggestive: $\kappa$ is literally a coordinate of the quartic family,
entering its two lowest coefficients as $e_3=e_1(1-\kappa)$ and
$e_4=(1-\kappa)^2$, so ``the quartic matches data'' and ``$\kappa=3/8$''
are two readings of one algebraic object evaluated at
$(I,\kappa)=(\tfrac34,\tfrac38)$. Given Theorem~\ref{thm:quartic}, three
formulations --- kernel, quartic, custodial coefficient --- are now built
on the single Casimir pair $(\tfrac34,2)$. Whether this is
over-determination or one coincidence wearing three coats is exactly the
question a trials-corrected analysis must answer; the equivalence results
of this paper are necessary bookkeeping input to that computation, since
they prevent double counting.

\section{Toward a geometric origin}
\label{sec:geometry}

Nothing above is a dynamical theory: \eqref{eq:kernel} is a spectral
ansatz. We now state what a derivation must produce, and why
Kaluza--Klein-type geometry is the natural habitat, without committing to a
specific internal manifold.

\begin{assumption}[Internal label]
\label{ass:su2i}
There is an internal $SU(2)_i$ --- an isometry of an internal space or a
custodial remnant --- under which the electroweak bosons carry
$(\gamma,W,Z)=(0,\tfrac12,1)$, and which is approximately conserved by the
dynamics that builds the tower.
\end{assumption}

\begin{assumption}[Pairing through the generators]
\label{ass:pairing}
The visible state $|v\rangle$ of the spin-$j$ multiplet couples to its
partner states through the internal generators $\Jvec$ themselves, with one
universal normalization: the coupling row vector is
$\langle p|\Jvec|v\rangle$ over partner states $|p\rangle$. Its squared
norm is
\[
\sum_p \bigl|\langle p|\Jvec|v\rangle\bigr|^2
=\langle v|\Jvec^{\,2}|v\rangle = j(j+1) = J
\]
for \emph{every} state of \emph{every} representation, so the two-block
reduction onto $\bigl(|v\rangle,\ \Jvec|v\rangle/\sqrt J\bigr)$ has
off-diagonal exactly $\sqrt J$, uniformly in $j$, integer or half-integer,
and independently of any basis choice. (A per-component version --- coupling
through a single ladder matrix element $\langle j,m{\pm}1|J_\pm|j,m\rangle$
--- reproduces $\sqrt J$ only for integer $j$ from the $m=0$ state and
fails for half-integer $j$; the operator-norm version above supersedes it.)
\end{assumption}

\begin{assumption}[Partner mass, sign, and nesting]
\label{ass:partner}
The partner multiplet carries diagonal mass-squared $J$ in the same units
--- the Casimir eigenvalue of an internal Laplacian, as on any coset space,
where the relevant spectra are universal multiples of $C_2$
\cite{salamstrathdee1982,duffnilssonpope1986} --- with the tachyonic sign
and the nested-decimation structure of
Proposition~\ref{prop:ladder}(iii), for which curvature-induced negative
vector-mode directions on distorted internal spaces
\cite{rdss1983,duffnilssonpope1986} and Wilson-line (Hosotani) dynamics
\cite{hosotani1983} are the known geometric sources.
\end{assumption}

Under Assumptions \ref{ass:su2i}--\ref{ass:partner} the spectrum
\eqref{eq:masses} follows. Assumption~\ref{ass:pairing} is the least
negotiable ingredient and also the most suggestive: the equality of the
coupling with $\sqrt{C_2}$ --- the corpus's ``coefficient lock'' --- holds
automatically when the mixing operator \emph{is} the generator, normalized
by the same constant that normalizes the Casimir. That is precisely the
structure of a gauge coupling to isometry currents, which is why we regard
Kaluza--Klein geometry
\cite{manton1979,forgacsmanton1980,witten1981,salamstrathdee1982} as the
natural habitat: any mechanism producing instead a mixing $g\Jvec$ with an
independent strength $g\ne1$ is incompatible with the observed five-digit
agreement of $m_W/m_Z$. What no calculation in the corpus or, to our
knowledge, in the literature has yet produced is the conjunction of all
three assumptions in one geometry --- with Remark~\ref{rem:notchain} now
adding a structural constraint: the geometry must realize a nested
decimation, and flat-chain constructions are excluded. The corpus's
classification-sweep program (enumerating coset geometries and testing
which, if any, produce the lock) is the right instrument, and its outcome
--- a derivation or a clean no-go --- will decide whether this section
describes physics or bookkeeping.

\section{Failure conditions}
\label{sec:falsification}

The construction is falsifiable on several independent fronts.
\begin{enumerate}
\item \emph{$m_W$ precision.} The tree-level on-shell prediction is
$m_W=80.3745(19)\GeV$ given $m_Z$. A settled world average more than
${\sim}2\sigma$ away kills the positive branch as an exact on-shell
statement; the 2024--2026 trend (CMS at $+1.4\sigma$, the 2026 snapshot at
$+1.5\sigma$) is mildly unfavorable and moving, while any confirmation of
the CDF~II region above $80.40\GeV$ excludes it outright.
\item \emph{Scheme stability.} If the one-loop matching program shows that
no consistent scheme assignment reproduces the algebraic number within
errors, the five-digit match is an on-shell accident.
\item \emph{The tower.} LEP~2 recasts above $193\GeV$ and FCC-ee either
find neutral states at $96.5/99.6/101.5/102.7\GeV$ with the predicted
weak-single/allowed-pair pattern, or exclude the tower up to
$106.58\GeV$, falsifying the extended assignment outright.
\item \emph{Negative branch.} Lemma~\ref{lem:incompatible} forces a
choice; independent precision on $v$, $m_t$, and the Higgs potential pins
both conventional scales, and the surviving identification then predicts
the other's offset ($1.8\%$ in ratio) with no free parameter.
\item \emph{Statistics.} A trials-corrected analysis over the full choice
space (Section~\ref{sec:provenance}) may return an unremarkable
probability --- the soft but decisive failure mode for the entire
enterprise.
\end{enumerate}

\section{Method, provenance, and trials accounting}
\label{sec:provenance}

This manuscript was produced by a language model working over a corpus of
research repositories, under a constraint of zero code execution: every
closed form was derived symbolically in the model's working notes, and
every decimal was obtained by hand arithmetic from those closed forms,
rounded at the last digit displayed. The constraint delimits the capability
being demonstrated: mathematical synthesis --- noticing that two
formulations are one (Theorem~\ref{thm:quartic}), that a truncation is
secretly a resummation with a spectral obstruction
(Proposition~\ref{prop:ladder}, Remark~\ref{rem:notchain}), that two
identifications are mutually inconsistent (Lemma~\ref{lem:incompatible})
--- rather than numerical search. An adversarial referee round by an
independent model materially improved the paper: it caught a false identity
in the first draft's irreducibility argument (since replaced by the orbit
argument in Theorem~\ref{thm:quartic}), forced the honest reframing of
Proposition~\ref{prop:ladder}, and corrected a mis-stated experimental
pull.

Provenance, stated exactly:
\begin{itemize}
\item \emph{Inherited:} the empirical observation and channel assignment
\eqref{eq:channels} \cite{devries2004,devriesrivero2005,rivero2006}; the
secular equation and its two-channel reading, with the mixing value
$0.22310132$ stated already in \cite{rivero2006}; the quartic
\eqref{eq:quartic}, the negative-branch identifications
\eqref{eq:negmatch}, and the $\kappa=3/8$ EFT result, all from the
unpublished corpus; the experimental inputs \cite{pdg2024,cmsw2024,cdf2022}.
\item \emph{Reproved from scratch:} the closed spectrum
(Section~\ref{sec:spectrum}); the convergence of the decimation recursion.
\item \emph{New here, to our knowledge:} the proof of
Theorem~\ref{thm:quartic} (the corpus had verified the coincidence
numerically, never algebraically); the identification of the quartic's
four roots with the four branch combinations and the two-parameter family
of Theorem~\ref{thm:family}, whose coefficients are polynomial in
$(I,\kappa)$ and which hands the statistics program its natural
hypothesis class; Corollary~\ref{cor:rho} with the degree-$8$ proof; the
flat-chain obstruction of Remark~\ref{rem:notchain}, including the
threshold $J=\tfrac94$; the operator-norm formulation of the coefficient
lock in Assumption~\ref{ass:pairing}, which repairs the half-integer case;
the branch sum rules in linking form (Proposition~\ref{prop:sumrules});
the no-go of Lemma~\ref{lem:incompatible} with its $1.8\%$ figure and the
compact reformulation $v=(1+\sqrt3)\,m_Z$ of the $J{=}2$ identification
\eqref{eq:vformula}; the label-conservation pattern of
Section~\ref{sec:selection} with the LEP~2 recast and the $193\GeV$ pair
threshold; and the observation that the biquadratic field makes $\swsq$
ruler-and-compass constructible.
\end{itemize}

\emph{Trials accounting.} Any significance claim for this construction must
price, at minimum, the following post-hoc choices: the kernel family
\eqref{eq:kernel} among comparably simple $2\times2$ textures; the map
$J=j(j+1)$ rather than another Casimir function; the positive-branch
selection; the channel assignment \eqref{eq:channels}; the anchor at $m_Z$;
the on-shell scheme; the extension of the tower and the association with
the $95$--$96\GeV$ hints; and the negative-branch matches
\eqref{eq:negmatch}. The corpus contains Monte Carlo look-elsewhere
machinery and quotes trials-corrected probabilities of order
$10^{-4}$--$10^{-6}$ for parts of this structure; those numbers are
unaudited here and no significance is claimed on their basis. What this
paper contributes to that future audit is subtractive and additive at once:
Theorem~\ref{thm:quartic} collapses two formulations into one (reducing the
apparent confirmation count), while Lemma~\ref{lem:incompatible} shows the
negative-branch sector currently \emph{fails} an exactness test --- both
facts that an honest significance computation must incorporate.

The honest prior on constructions of this genre is unfavorable: the history
of electroweak numerology is a graveyard, and the Koide relation itself,
after forty years of real theoretical effort
\cite{koide1990,sumino2009,riverogsponer2005}, remains a suggestive orphan.
The construction studied here has three properties that graveyard occupants
typically lack: a five-digit match with one input and no continuous
parameters; an internal algebraic rigidity that generates new exact
statements when probed (Sections \ref{sec:theorem}, \ref{sec:negative});
and a finite list of near-term experimental executioners
(Section~\ref{sec:falsification}). We submit that this is the correct
standard for work of this kind: a well-posed falsification program.

\section*{Acknowledgments}
The corpus of repositories synthesized here is the work of A.~Rivero and
collaborating AI systems; the Koide-mechanism reference notes and
literature survey in that corpus were directly useful. The bibliography was
verified against arXiv and the published record by a Claude Sonnet agent;
the adversarial review was performed by GPT-5.5 (OpenAI) via the Codex CLI.
Both reports are archived with the source.

\begin{thebibliography}{99}

\bibitem{devries2004} H.~de~Vries,
online post on the electroweak mixing-angle coincidence,
Physics Forums, thread ``All the lepton masses from $G$, $\pi$, $e$,''
post \#44 (November 2004),
\url{https://www.physicsforums.com/showpost.php?p=382642&postcount=44}.
Earliest public statement of the relation; documented and dated in
\cite{rivero2006}, Sec.~3.

\bibitem{devriesrivero2005} H.~de~Vries and A.~Rivero,
``Evidence for radiative generation of lepton masses,''
arXiv:hep-ph/0503104.

\bibitem{rivero2006} A.~Rivero,
``Mass terms to break susy-like degeneration,''
arXiv:hep-ph/0606171. States the Casimir-operator construction whose
positive eigenvalues give $s^2_{dV}=1-m^2_{1/2,+}/m^2_{1,+}=0.22310132$.

\bibitem{koide1982} Y.~Koide,
``Fermion-boson two-body model of quarks and leptons and Cabibbo mixing,''
Lett.\ Nuovo Cim.\ {\bf 34} (1982) 201. See also
``A fermion-boson composite model of quarks and leptons,''
Phys.\ Lett.\ B {\bf 120} (1983) 161, which first states the $K=2/3$
formula in its standard form.

\bibitem{koide1983prd} Y.~Koide,
``New view of quark and lepton mass hierarchy,''
Phys.\ Rev.\ D {\bf 28} (1983) 252 [Erratum: Phys.\ Rev.\ D {\bf 29}
(1984) 1544].

\bibitem{koide1990} Y.~Koide,
``Charged lepton mass sum rule from U(3)-family Higgs potential model,''
Mod.\ Phys.\ Lett.\ A {\bf 5} (1990) 2319.

\bibitem{koide2005} Y.~Koide,
``Challenge to the mystery of the charged lepton mass formula,''
arXiv:hep-ph/0506247.

\bibitem{koide2018} Y.~Koide,
``What physics does the charged lepton mass relation tell us?,''
arXiv:1809.00425 [hep-ph].

\bibitem{foot1994} R.~Foot,
``A note on Koide's lepton mass relation,''
arXiv:hep-ph/9402242.

\bibitem{sumino2009} Y.~Sumino,
``Family gauge symmetry and Koide's mass formula,''
Phys.\ Lett.\ B {\bf 671} (2009) 477, arXiv:0812.2090 [hep-ph]. Companion
full model: JHEP {\bf 0905} (2009) 075, arXiv:0812.2103.

\bibitem{riverogsponer2005} A.~Rivero and A.~Gsponer,
``The strange formula of Dr.\ Koide,''
arXiv:hep-ph/0505220.

\bibitem{feshbach1958} H.~Feshbach,
``Unified theory of nuclear reactions,''
Annals Phys.\ {\bf 5} (1958) 357.

\bibitem{manton1979} N.~S.~Manton,
``A new six-dimensional approach to the Weinberg-Salam model,''
Nucl.\ Phys.\ B {\bf 158} (1979) 141.

\bibitem{forgacsmanton1980} P.~Forgacs and N.~S.~Manton,
``Space-time symmetries in gauge theories,''
Commun.\ Math.\ Phys.\ {\bf 72} (1980) 15.

\bibitem{witten1981} E.~Witten,
``Search for a realistic Kaluza-Klein theory,''
Nucl.\ Phys.\ B {\bf 186} (1981) 412.

\bibitem{salamstrathdee1982} A.~Salam and J.~Strathdee,
``On Kaluza-Klein theory,''
Annals Phys.\ {\bf 141} (1982) 316.

\bibitem{rdss1983} S.~Randjbar-Daemi, A.~Salam and J.~Strathdee,
``Instability of higher dimensional Yang-Mills systems,''
Phys.\ Lett.\ B {\bf 124} (1983) 345 [Erratum: Phys.\ Lett.\ B {\bf 144}
(1984) 455].

\bibitem{hosotani1983} Y.~Hosotani,
``Dynamical mass generation by compact extra dimensions,''
Phys.\ Lett.\ B {\bf 126} (1983) 309.

\bibitem{duffnilssonpope1986} M.~J.~Duff, B.~E.~W.~Nilsson and C.~N.~Pope,
``Kaluza-Klein supergravity,''
Phys.\ Rept.\ {\bf 130} (1986) 1.

\bibitem{pdg2024} S.~Navas {\it et al.} [Particle Data Group],
``Review of particle physics,''
Phys.\ Rev.\ D {\bf 110} (2024) 030001.

\bibitem{cmsw2024} [CMS Collaboration],
``Measurement of the W boson mass in proton-proton collisions at
$\sqrt s=13$ TeV,'' arXiv:2412.13872 [hep-ex].

\bibitem{cdf2022} T.~Aaltonen {\it et al.} [CDF Collaboration],
``High-precision measurement of the W boson mass with the CDF II
detector,'' Science {\bf 376} (2022) 170.

\bibitem{cms95} [CMS Collaboration],
``Search for a standard model-like Higgs boson in the mass range between
70 and 110 GeV in the diphoton final state in proton-proton collisions at
$\sqrt s=13$ TeV,'' Phys.\ Lett.\ B {\bf 860} (2024) 139067,
arXiv:2405.18149 [hep-ex].

\bibitem{atlas95} [ATLAS Collaboration],
``Search for diphoton resonances in the 66 to 110 GeV mass range using
$pp$ collisions at $\sqrt s=13$ TeV with the ATLAS detector,''
JHEP {\bf 01} (2025) 053, arXiv:2407.07546 [hep-ex].

\end{thebibliography}

\end{document}
